% ---- pg-030 （书页 18） ----

\begin{mvfigure}
\mvinc{pg-030}{0.18513,0.13458,0.80213,0.41280}
\phantomsection\label{fig:1.12}\mvcap{图 1.12\quad 对木材硬度数据的三维透视}
\end{mvfigure}

像\hyperref[fig:1.12]{图 1.12} 这样的图使我们可以迅速地识别出那些与其他数据不一致的观测值和可能会严重影响到基于标准数据生成模型所得推论的观测值.

\mvsub{生长曲线图}

如果在每年生日的时候都测一下小孩的身高,然后将点画出来,连接成图.这就是一个典型的生长曲线图.一般来说,如果测量值随着时间有上升、下降或者既有上升又有下降的情况,我们就在相同的单位和对象下,可以重复测量该值,得到生长曲线图.

\mvbold{例 1.10} (生长曲线阵列)

阿拉斯加渔业和狩猎部门为了保持灰熊健康的种群量,对灰熊实施监测.熊被催眠针催眠然后挂在三脚架上测量它们的体重,并用钢卷尺测量它们的身长.\hyperref[tab:1.4]{表 1.4} 中给出了 7 只母熊在 2,3,4 和 5 岁时的体重(简记为 Wt,单位: 公斤)和身长(简记为 Lngth,单位: 厘米)的数据.

\begin{mvtable}
\begin{tabular}{ccccccccc}
\toprule
熊序号 & Wt2 & Wt3 & Wt4 & Wt5 & Lngth2 & Lngth3 & Lngth4 & Lngth5 \\
\midrule
1 & 48 & 59 & 95 & 82 & 141 & 157 & 168 & 183 \\
2 & 59 & 68 & 102 & 102 & 140 & 168 & 174 & 170 \\
3 & 61 & 77 & 93 & 107 & 145 & 162 & 172 & 177 \\
4 & 54 & 43 & 104 & 104 & 146 & 159 & 176 & 171 \\
5 & 100 & 145 & 185 & 247 & 150 & 158 & 168 & 175 \\
6 & 68 & 82 & 95 & 118 & 142 & 140 & 178 & 189 \\
7 & 68 & 95 & 109 & 111 & 139 & 171 & 176 & 175 \\
\bottomrule
\end{tabular}
\phantomsection\label{tab:1.4}\mvtabcap{表 1.4\quad 母熊数据}
\end{mvtable}

{\zihao{6} 资料来源: 数据经 H. Roberts 许可使用.}
